\chapter{总结与展望}\label{chap:conclusion}

\section{本文工作总结}

分布式集群是支撑现代云计算与大数据处理的核心基础设施，其稳定性直接关系到上层业务的连续性。针对集群异常检测中现有方法未能有效利用多节点行为相似性这一重要先验知识的问题，本文提出了基于多元时序曲线相似性的分布式集群异常检测方法。本文的主要工作和贡献总结如下。

第一，本文提出了一种基于多节点曲线相似性的异常检测方法论。该方法突破了传统异常检测聚焦于时间与指标二维分析的框架，创新性地将集群节点这一空间维度系统性地引入模型，构建了时间、指标与节点的三维分析范式。通过利用分布式集群负载均衡机制所保证的节点行为相似性，将隐性的运维经验转化为显式的可计算先验知识，有效规避了因定义绝对正常模式而产生的误报问题。

第二，针对单指标多节点场景，本文基于联通大数据生产环境构建了高质量的时序数据集，并分别采用欧氏距离、自编码器嵌入和DBSCAN聚类三种方法度量节点曲线间的相似性。针对单一方法阈值敏感的问题，设计了基于逻辑与策略的集成学习框架，并引入改进的包容性SPOT算法实现阈值的动态自适应调整。实验结果表明，所提方法在F1分数上显著优于现有的多元时序异常检测算法。

第三，针对多指标多节点场景，本文提出了数值与形状双视角融合的异常检测框架。在数值视角方面，采用基于欧氏距离的多指标加权融合方法量化节点间的数值差异。在形状视角方面，提出了改进的SAX符号化表示算法，通过保留原始时间分辨率和采用严格距离计算规则来提升对形状变化的敏感性。在GPU集群故障数据集上的实验表明，融合方法的故障检测准确率达到83.8\%，显著优于单一视角方法和传统基线方法。

第四，本文设计并实现了基于曲线相似性的异常检测系统，实现了从数据采集、预处理、异常检测到结果展示的完整流程。系统已在联通大数据生产底座平台上进行部署验证，成功识别出多起节点级异常事件，验证了方法在实际生产环境中的可行性和有效性。

\section{未来研究展望}

虽然本文在基于多节点曲线相似性的异常检测方面取得了一定进展，但仍有若干值得深入研究的方向。

第一，更复杂的节点关系建模。本文假设同构集群内各节点地位平等，但实际系统中节点可能存在主从关系、依赖关系等复杂结构。未来可以引入图神经网络等方法显式建模节点间的拓扑关系，提升对复杂系统结构的刻画能力。

第二，异构集群的异常检测。本文主要关注同构集群中节点行为的相似性，但在异构集群中，不同类型节点的正常行为模式可能存在差异。未来可以研究如何在异构环境中识别可比较的节点组，并设计适应性更强的检测方法。

第三，在线增量学习能力。当前方法主要采用批量处理模式，对于持续运行的系统，需要能够在线更新模型以适应系统行为的长期演变。未来可以研究增量学习方法，使模型能够在保持对历史模式记忆的同时适应新的正常行为模式。

第四，异常根因分析。当前方法主要关注异常检测和故障节点定位，对于异常的具体原因缺乏深入分析。未来可以结合日志分析、配置变更追踪等信息，构建从异常检测到根因定位的完整链路，进一步提升系统的智能运维能力。

第五，跨集群知识迁移。不同集群的监控数据虽然在绝对值上可能存在差异，但其中蕴含的异常模式可能具有相似性。未来可以研究如何在多个集群之间进行知识迁移和经验共享，提升检测方法在新场景中的泛化能力。
